\section{Entrevista al Gerente de Cafetalino}\label{anexo-entrevista-al-gerente-de-cafetalino}

\textbf{GUÍA PARA LA ENTREVISTA AL GERENTE PROPIETARIO DE CAFETALINO}

El objetivo de la entrevista es conocer y comprender la empresa y sus procedimientos a fin de identificar problemas que podrían encararse con la ayuda de sistemas basados en inteligencia artificial.

\begin{enumerate}
\def\labelenumi{\arabic{enumi}.}
\item
  ¿Tiene identificado algún problema recurrente, proceso pesado, o algo que le gustaría optimizar? Por ejemplo, predecir qué meses se venderá más o automatizar alguna tarea.
\item
  ¿Guarda algún tipo de datos o registros de sus operaciones?
\item
  ¿Considera que una aplicación de inteligencia artificial aportaría valor a sus operaciones?
\item
  ¿Estaría dispuesto a compartirla para implementar soluciones de inteligencia artificial con el objetivo de mejorar sus operaciones?
\item
  ¿Alguna vez intento instalar sensores para hacer seguimiento al consumo de sus insumos?
\item
  ¿Alguna vez ha tenido algún incidente con sus clientes? ¿De qué tipo?
\item
  ¿Cómo ha resuelto estos problemas?
\item
  ¿Cuál es el procedimiento de recarga de sus máquinas dispensadoras?
\item
  ¿Alguna vez ha tenido algún problema con este proceso? ¿De qué tipo?
\end{enumerate}

\textbf{Desarrollo de la Entrevista al Gerente Propietario de Cafetalino}

El objetivo de la entrevista es conocer y comprender la empresa y sus procedimientos a fin de identificar problemas que podrían encararse con la ayuda de sistemas basados en inteligencia artificial.

\begin{enumerate}
\def\labelenumi{\arabic{enumi}.}
\item
  ¿Tiene identificado algún problema recurrente, proceso pesado, o algo que le gustaría optimizar? Por ejemplo, predecir qué meses se venderá más o automatizar alguna tarea.
\end{enumerate}

\begin{quote}
\emph{Tenemos muchas cosas que queremos incorporar y mejorar a lo que actualmente hacemos:}
\end{quote}

\begin{enumerate}
\def\labelenumi{\arabic{enumi}.}
\item
  \emph{Enviar notificaciones a los clientes de que su tarjeta esta con crédito insuficiente y que les sugiera recargar.}
\item
  \emph{Nuestro sistema lleva registro de todas las operaciones con tarjeta. Nos gustaría que también registre las operaciones con monedas y que esto este diferenciado: Operaciones con tarjeta y operaciones con monedas.}
\item
  \emph{Quisiéramos abrir nuestro sistema a los clientes mediante una app y/o una web para consultas de saldo, recargas de tarjeta u otros beneficios}
\item
  \emph{Nos gustaría poder conocer, de todas las recargas a las tarjetas, cuantas de ellas ya se han consumido y cuánto es el saldo, que se constituye en producto por entregar, que ha sido pagado por adelantado. Así hay muchas cosas, necesitamos muchas cosas más, hay muchas ideas.}
\item
  \emph{Por ejemplo, un sistema de inventario de insumos, que son distribuidos a las maquinas, estos se convierten en productos y desde las maquinas son comercializados. Que me devuelva saldos para verificarlos físicamente o que pueda tener una tendencia para pronosticar los requerimientos de recarga de insumos a las maquinas.}
\end{enumerate}

\begin{quote}
Esta quinta idea me parece interesante y viable para darle una solución con inteligencia artificial, que es lo que me gustaría aplicar.
\end{quote}

\begin{enumerate}
\def\labelenumi{\arabic{enumi}.}
\setcounter{enumi}{1}
\item
  ¿Guarda algún tipo de datos o registros de sus operaciones?
\end{enumerate}

\begin{quote}
\emph{Si, tenemos datos de toda la actividad de Cafetalino desde el 2018, de hecho, también se podría estimar la cantidad de insumos que hemos usado en función a las ventas.}

\emph{Hay un registro de todas las operaciones con tarjeta. Nos gustaría que también registre las operaciones con monedas y que esto este diferenciado: Operaciones con tarjeta y operaciones con monedas.}
\end{quote}

\begin{enumerate}
\def\labelenumi{\arabic{enumi}.}
\setcounter{enumi}{2}
\item
  ¿Considera que una aplicación de inteligencia artificial aportaría valor a sus operaciones?
\end{enumerate}

\begin{quote}
\emph{Por supuesto, además sería una forma de modernizar nuestras operaciones, de estar a la vanguardia}
\end{quote}

\begin{enumerate}
\def\labelenumi{\arabic{enumi}.}
\setcounter{enumi}{3}
\item
  ¿Estaría dispuesto a compartir su información para implementar soluciones de inteligencia artificial con el objetivo de mejorar sus operaciones?
\end{enumerate}

\begin{quote}
\emph{Mmm, si, para mejorar nuestras operaciones por supuesto que sí, pero la condición es manejar con mucha discreción, no quisiéramos que, de ninguna manera, se filtre la información y cosas que nos ha costado aprender sean aprovechadas por otros para hacernos la competencia. De hecho, tendríamos que hacer un contrato de confidencialidad.}

Por supuesto, no hay de que preocuparse pues se manejaría con mucha profesionalidad.
\end{quote}

\begin{enumerate}
\def\labelenumi{\arabic{enumi}.}
\setcounter{enumi}{4}
\item
  ¿Alguna vez intento instalar sensores para hacer seguimiento al consumo de sus insumos?
\end{enumerate}

\begin{quote}
\emph{Si, en una primera instancia se realizó una prueba piloto y se instalaron sensores en algunas máquinas para saber si era necesario recargar vasos, y que el sistema nos alerte, vía mensaje, cuando quedaba una pequeña cantidad. Sin embargo, se tuvieron problemas de conectividad y no prospero la idea, además se limitaba solo a los vasos y no a los demás insumos.}

\emph{Se exploro la posibilidad de instalar más sensores, uno para cada uno de insumos, pero el presupuesto inicial era muy alto, no había la certeza de que funcionara, pues los problemas de conectividad persistían, además había limitaciones técnicas especialmente con los insumos sólidos, yyy también con los líquidos.}
\end{quote}

\begin{enumerate}
\def\labelenumi{\arabic{enumi}.}
\setcounter{enumi}{5}
\item
  ¿Alguna vez ha tenido algún incidente con sus clientes? ¿De qué tipo?
\end{enumerate}

\begin{quote}
\emph{Si. Alguna vez cuando una maquina se ha quedado sin insumos o quedó muy poca cantidad de alguno de ellos, ocurrió que: o no entrego el producto o este no tenía la calidad a la que el cliente estaba acostumbrado. Esto genera, de forma inmediata sentimientos de frustración en el cliente, decepción e impotencia, no saben a quién reclamar, y una sensación de pérdida o robo, pierden la confianza y muchos ya no se animan a volver a interactuar con la máquina.}

\emph{En algunos casos estos sentimientos escalaron en el nivel de violencia, algunos usuarios con menor tolerancia, tuvieron comportamientos agresivos, golpearon los aparatos, y en ocasiones provocaron daños a nuestras maquinas; en otras ocasiones llamaron al número de contacto y nos maltrataron, nos acusaron de robo y no aceptaron nuestras disculpas, ni nos dejaron darles ninguna explicación, hasta nos amenazaron con procesos judiciales, nos difamaron en las redes sociales y en su entorno de trabajo. Alguna vez hasta tuvimos ir a su entorno de trabajo, les explicamos lo sucedido y a modo de disculpa les ofrecimos bebidas a todos los trabajadores.}

\emph{En un caso en particular, recibimos la denuncia de sus compañeros de trabajo, que el sujeto provocaba el fallo de forma intencionada, introduciendo otros objetos o monedas de otros países y que luego protestaba en voz alta. Bueno, también nos comentaron que esa era su forma cotidiana de actuar y hasta tuvimos que retirar nuestro equipo de esa ubicación para evitar el daño, a los equipos y a la imagen de la empresa.}
\end{quote}

\begin{enumerate}
\def\labelenumi{\arabic{enumi}.}
\setcounter{enumi}{6}
\item
  ¿Cómo ha resuelto estos problemas?
\end{enumerate}

\begin{quote}
\emph{Bueno, en estas situaciones, los que nos conocen nos llamaban y nos comunicaban lo que sucedió; nosotros le devolvíamos su dinero y le dábamos de cortesía la bebida, cuando era posible.}

\emph{Luego pusimos un número de teléfono de contacto en todas nuestras ubicaciones, para que cualquier cliente, aunque no nos conozca, pueda reportar el suceso y procedíamos de la misma manera, devolverle su dinero y si era posible darle su bebida de cortesía}

\emph{Sin embargo, como obviamente transcurría un tiempo para que podamos ir a recargar la maquina y ponerla nuevamente en servicio y no siempre era posible darle la bebida de cortesía. Con la devolución del dinero, no hubo problema, inicialmente hablamos con vecinos de nuestras ubicaciones que nos hacían el favor de devolver el dinero y luego nosotros les reponíamos, pero ahora con la facilidad de las trasferencias por Qr eso se ha facilitado enormemente}.
\end{quote}

\begin{enumerate}
\def\labelenumi{\arabic{enumi}.}
\setcounter{enumi}{7}
\item
  ¿Cuál es el procedimiento de recarga de sus máquinas dispensadoras?
\end{enumerate}

\begin{quote}
\emph{Al principio cuando teníamos pocas maquinas, y nos estábamos haciendo conocer, hacíamos el recorrido por una ruta fija, revisando todas nuestras máquinas y recargando en función a lo consumido. Esto lo hacíamos una vez por semana, una vez cada 2 semanas, una vez cada 10 días, fuimos ajustando la frecuencia en base a la cantidad consumida.}

\emph{Después conforme nos fuimos haciendo conocer, aumentamos la frecuencia de recarga y la cantidad de máquinas, luego agrupamos las maquinas en función a su ubicación y aumentamos el número de rutas. Actualmente, revisamos dos rutas cada día, y tenemos personal que se ocupa de ello.}
\end{quote}

\begin{enumerate}
\def\labelenumi{\arabic{enumi}.}
\setcounter{enumi}{8}
\item
  ¿Alguna vez ha tenido algún problema con este proceso? ¿De qué tipo?
\end{enumerate}

\begin{quote}
\emph{Bueno, si se le puede llamar problema, a veces ocurre que, al revisar una máquina, no ha tenido el consumo esperado y no requiere recarga de insumos.}

\emph{En otras ocasiones, nos hemos visto en la necesidad de modificar nuestra ruta porque se nos reportó que una máquina, de la misma ruta se quedó sin insumos y dejo de prestar el servicio, de manera que es necesario priorizarla, o cuando se trata de una máquina de otra ruta, es necesario priorizarla e incluirla en la secuencia. Esto nos obliga a llevar una cantidad extra de insumos. En otras ocasiones el reporte de una maquina fuera de servicio, llegó al terminar el recorrido y se tuvo que volver a realizar la rutina por una sola máquina: preparar paquetes de recarga de insumos, dirigirnos a la ubicación lo más rápido posible, revisar la máquina, hacerle un mantenimiento rutinario, reabastecerla y ponerla en servicio.}

¿Tiene idea de cuántas veces ocurren estos eventos a la semana, al mes o al año?

\emph{Maquinas que no tienen necesidad de recarga aproximadamente 1 de cada 10.}

\emph{Eventos de modificación de la ruta de reabastecimiento, 10 al año.}

Se me ocurre una aplicación que permita que procesar el historial de recargas (2018-2026) para predecir el desabastecimiento y planificar automáticamente rutas dinámicas óptimas.

\emph{Perfecto, me parece interesante la idea.}
\end{quote}
